\documentclass[11pt]{article}
\usepackage[margin=0.95in]{geometry}
\usepackage{amsmath,amssymb,mathtools,bm}
\usepackage[hidelinks]{hyperref}
\usepackage[capitalise]{cleveref}

\newcommand{\wt}{\widetilde}
\newcommand{\dd}{\mathrm{d}}
\newcommand{\ii}{\mathrm{i}}
\newcommand{\abs}[1]{\left|#1\right|}
\newcommand{\one}{\mathbf 1}
\DeclareMathOperator{\Res}{Res}

\title{Audited Corrected Analytical Derivation\\
Nearest-neighbour ring walk with one absorbing site and one shortcut}
\author{}
\date{\today}

\begin{document}
\maketitle

\tableofcontents

\section{Purpose and final conclusions}

This document is a detailed audit and correction of the analytical calculation for a
nearest-neighbour discrete-time random walk on a ring with one absorbing site and one directed
shortcut into the absorbing site.  The main conclusion is that the absorbing-ring part of the
calculation can be made correct after several index, sign, and orthogonality corrections, but the
shortcut-to-target part must be changed more substantially.

The key point is the following.  If a shortcut from $u$ to the absorbing target $v$ replaces a
fraction $\beta$ of the stay-put probability $1-q$, then in the transient state space the shortcut is
an extra killing channel at $u$.  Writing
\begin{equation}
  \lambda=\beta(1-q),\qquad 0\leq \beta\leq 1,
\end{equation}
the corrected resolvent is
\begin{equation}
  \boxed{
  \wt S_{n_0}(n,z)=\wt W_{n_0}(n,z)
  -\frac{z\lambda\,\wt W_{n_0}(u,z)\wt W_u(n,z)}
  {1+z\lambda\,\wt W_u(u,z)}.}
  \label{eq:main-corrected-resolvent}
\end{equation}
The minus sign and the plus sign in the denominator are both forced by the stochastic
interpretation: the shortcut moves probability mass from the transient site $u$ to the absorbing
site $v$.

Consequently, the first-passage generating function is
\begin{equation}
  \boxed{
  \wt F_{n_0}(v,z)=\wt F^{(0)}_{n_0}(v,z)+
  \frac{z\lambda\,\wt W_{n_0}(u,z)\bigl[1-\wt F^{(0)}_u(v,z)\bigr]}
  {1+z\lambda\,\wt W_u(u,z)}.}
  \label{eq:main-corrected-F}
\end{equation}
Here the superscript $(0)$ means ``without the shortcut but with the absorbing target present.''
This formula also gives a useful monotonicity check: the shortcut contribution to
$\wt F_{n_0}(v,z)$ has non-negative coefficients for $z\in[0,1)$, so adding a shortcut to the
target cannot decrease the probability of having hit the target by any fixed time.

In the symmetric case
\begin{equation}
  N=2L,
  \qquad v=0,
  \qquad u=L,
\end{equation}
and for an initial site at distance $\rho$ from $v$ along either half-ring,
$0\leq\rho\leq L$, the corrected first-passage generating function reduces to the closed form
\begin{equation}
  \boxed{
  \wt F_{\rho}(z)=
  \frac{aT_{L-\rho}(y)+U_{\rho-1}(y)+U_{L-\rho-1}(y)}
  {aT_L(y)+U_{L-1}(y)},
  \qquad
  a=\frac{q}{\beta(1-q)},
  \qquad
  y=1+\frac{1}{q}\left(\frac1z-1\right).}
  \label{eq:main-symmetric-F}
\end{equation}
We use the convention $U_{-1}=0$.  Formula \eqref{eq:main-symmetric-F} is the cleanest corrected
analytical answer.

The poles are the roots of
\begin{equation}
  \boxed{aT_L(y_s)+U_{L-1}(y_s)=0,
  \qquad y_s=1+\frac{s-1}{q}.}
  \label{eq:main-pole-equation-y}
\end{equation}
Equivalently, if $y_s=\cos\theta$, the roots satisfy
\begin{equation}
  \boxed{\tan(L\theta)=-a\sin\theta.}
  \label{eq:main-pole-equation-theta}
\end{equation}
The largest positive root $s_1$ lies in
\begin{equation}
  \boxed{\gamma_1<s_1<\alpha_1,}
  \qquad
  \gamma_1=1-q+q\cos\frac{2\pi}{N},
  \qquad
  \alpha_1=1-q+q\cos\frac{\pi}{N}.
  \label{eq:s1-between}
\end{equation}
Thus the shortcut-to-target makes the long-time tail decay faster than the no-shortcut
absorbing-ring tail.  The original condition
\begin{equation}
  \beta\leq \frac{2q}{(1-q)N}
\end{equation}
is not a valid corrected condition for bimodality or pole dominance.  The corrected long-time
decay is
\begin{equation}
  \boxed{F_\rho(t)\sim B_{\rho 1}s_1^{t-1},
  \qquad t\to\infty,}
  \label{eq:main-tail}
\end{equation}
not $t\alpha_1^t$.  For weak shortcuts,
\begin{equation}
  \boxed{
  s_1=\alpha_1-\frac{2\beta(1-q)}{N}+O\bigl([\beta(1-q)]^2\bigr),}
  \label{eq:weak-shortcut-s1}
\end{equation}
and therefore
\begin{equation}
  \boxed{
  1-s_1=1-\alpha_1+\frac{2\beta(1-q)}{N}+O\bigl([\beta(1-q)]^2\bigr).}
  \label{eq:weak-shortcut-gap}
\end{equation}
For large $N$ this becomes
\begin{equation}
  \boxed{
  1-s_1\sim \frac{q\pi^2}{2N^2}+\frac{2\beta(1-q)}{N}.}
  \label{eq:large-N-gap}
\end{equation}
Equation \eqref{eq:large-N-gap} is the most direct corrected decay-dependence statement for the
right tail of the second peak, if a second peak is present.

\section{Model and notation}
\label{sec:model}

The ring sites are $0,1,\ldots,N-1$, with all site labels understood modulo $N$.  The unmodified
nearest-neighbour transition rule is
\begin{equation}
  X_{t+1}=\begin{cases}
  X_t, & \text{with probability }1-q,\\
  X_t+1, & \text{with probability }q/2,\\
  X_t-1, & \text{with probability }q/2,
  \end{cases}
  \qquad 0<q<1.
\end{equation}
The absorbing target is denoted by $m$ or, in the shortcut section, by $v$.  If $m$ is absorbing,
then once the walker reaches $m$ it is removed from the transient dynamics.  The killed propagator
is
\begin{equation}
  W_{n_0}(n,t)=\Pr\{X_t=n,\; X_0=n_0,\; X_s\neq m\text{ for all }0\leq s\leq t\},
\end{equation}
for transient sites $n,n_0\neq m$.  The first-passage probability to $m$ is
\begin{equation}
  \mathcal F_{n_0}(m,t)=\Pr\{X_t=m,\; X_s\neq m\text{ for }0\leq s<t\},
  \qquad t\geq 1,
\end{equation}
with $\mathcal F_{n_0}(m,0)=0$ when $n_0\neq m$.

Generating functions are defined by
\begin{equation}
  \wt W_{n_0}(n,z)=\sum_{t=0}^{\infty}W_{n_0}(n,t)z^t,
  \qquad
  \wt F_{n_0}(m,z)=\sum_{t=0}^{\infty}\mathcal F_{n_0}(m,t)z^t.
\end{equation}
The frequently used change of variable is
\begin{equation}
  \boxed{y=y(z)=1+\frac1q\left(\frac1z-1\right).}
  \label{eq:y-definition}
\end{equation}
It is useful because
\begin{equation}
  1-z\bigl[1-q+q\cos\theta\bigr]
  =zq\bigl[y-\cos\theta\bigr].
  \label{eq:denominator-y}
\end{equation}

When $N=2L$ is even, the odd-sector eigenvalues are
\begin{equation}
  \boxed{\alpha_\ell=1-q+q\cos\frac{(2\ell-1)\pi}{N},
  \qquad \ell=1,\ldots,L,}
  \label{eq:alpha-definition}
\end{equation}
and the even-sector ring eigenvalues are
\begin{equation}
  \boxed{\gamma_r=1-q+q\cos\frac{2\pi r}{N},
  \qquad r=1,\ldots,N-1.}
  \label{eq:gamma-definition}
\end{equation}

\section{Finite-sum identities with details}
\label{sec:finite-identities}

Let
\begin{equation}
  \theta_j=\frac{\pi j}{N},
  \qquad
  c_j=\cos\theta_j,
  \qquad
  A_j(d)=\bigl[1-(-1)^j\bigr]\sin(d\theta_j)\sin\theta_j,
\end{equation}
where $0<d<N$ is an integer.  The central identity is
\begin{equation}
  \boxed{
  \frac1N\sum_{j=1}^{N-1}\frac{A_j(d)}{y-c_j}
  =\frac{U_{N-d-1}(y)+U_{d-1}(y)}{U_{N-1}(y)}.}
  \label{eq:central-identity}
\end{equation}
Here $T_n$ and $U_n$ are Chebyshev polynomials of the first and second kind, with
\begin{equation}
  T_n(\cos\theta)=\cos(n\theta),
  \qquad
  U_n(\cos\theta)=\frac{\sin((n+1)\theta)}{\sin\theta},
  \qquad
  U_{-1}(y)=0.
\end{equation}

\subsection{Proof of the central identity}

The denominator $U_{N-1}(y)$ has the simple roots
\begin{equation}
  y=c_j=\cos\frac{\pi j}{N},
  \qquad j=1,\ldots,N-1.
\end{equation}
Indeed,
\begin{equation}
  U_{N-1}(c_j)=\frac{\sin(N\theta_j)}{\sin\theta_j}
  =\frac{\sin(\pi j)}{\sin\theta_j}=0.
\end{equation}
We compute the residue of the right-hand side of \eqref{eq:central-identity} at $y=c_j$.
The numerator is
\begin{align}
  U_{N-d-1}(c_j)+U_{d-1}(c_j)
  &=\frac{\sin((N-d)\theta_j)+\sin(d\theta_j)}{\sin\theta_j} \notag\\
  &=\frac{\sin(N\theta_j-d\theta_j)+\sin(d\theta_j)}{\sin\theta_j}.\label{eq:num-res-step1}
\end{align}
Since $N\theta_j=\pi j$,
\begin{align}
  \sin(N\theta_j-d\theta_j)
  &=\sin(\pi j)\cos(d\theta_j)-\cos(\pi j)\sin(d\theta_j)\notag\\
  &=0-(-1)^j\sin(d\theta_j)
  =-(-1)^j\sin(d\theta_j).
\end{align}
Substituting this into \eqref{eq:num-res-step1} gives
\begin{equation}
  U_{N-d-1}(c_j)+U_{d-1}(c_j)
  =\frac{\bigl[1-(-1)^j\bigr]\sin(d\theta_j)}{\sin\theta_j}.
  \label{eq:numerator-at-root}
\end{equation}
Next, differentiate $U_{N-1}(\cos\theta)=\sin(N\theta)/\sin\theta$ with respect to $y$.
Because $y=\cos\theta$, $\dd y/\dd\theta=-\sin\theta$.  Thus
\begin{equation}
  \frac{\dd}{\dd y}U_{N-1}(\cos\theta)
  =-\frac{1}{\sin\theta}\frac{\dd}{\dd\theta}
  \left[\frac{\sin(N\theta)}{\sin\theta}\right].
\end{equation}
At $\theta=\theta_j$, $\sin(N\theta_j)=0$ and $\cos(N\theta_j)=(-1)^j$, so
\begin{align}
  U_{N-1}'(c_j)
  &=-\frac{1}{\sin\theta_j}
  \left[\frac{N\cos(N\theta_j)\sin\theta_j-
  \sin(N\theta_j)\cos\theta_j}{\sin^2\theta_j}\right]\notag\\
  &=-\frac{1}{\sin\theta_j}\left[\frac{N(-1)^j\sin\theta_j}{\sin^2\theta_j}\right]
  =\frac{N(-1)^{j+1}}{\sin^2\theta_j}.
  \label{eq:U-prime-root}
\end{align}
Therefore the residue of the right-hand side at $y=c_j$ is
\begin{align}
  \Res_{y=c_j}\frac{U_{N-d-1}(y)+U_{d-1}(y)}{U_{N-1}(y)}
  &=\frac{U_{N-d-1}(c_j)+U_{d-1}(c_j)}{U_{N-1}'(c_j)}\notag\\
  &=\frac{\bigl[1-(-1)^j\bigr]\sin(d\theta_j)}{\sin\theta_j}
  \frac{\sin^2\theta_j}{N(-1)^{j+1}}.\label{eq:residue-rhs-before}
\end{align}
If $j$ is even, the factor $1-(-1)^j$ is zero.  If $j$ is odd, then
$(-1)^{j+1}=1$.  Hence in all cases
\begin{equation}
  \Res_{y=c_j}\frac{U_{N-d-1}(y)+U_{d-1}(y)}{U_{N-1}(y)}
  =\frac{1}{N}\bigl[1-(-1)^j\bigr]\sin(d\theta_j)\sin\theta_j
  =\frac{A_j(d)}{N}.
  \label{eq:residue-rhs-final}
\end{equation}
The left-hand side of \eqref{eq:central-identity} has exactly the same simple pole at $y=c_j$
and exactly the same residue $A_j(d)/N$.  The difference between the two sides has no finite
poles.  Both sides tend to zero as $y\to\infty$.  Therefore the difference is an entire rational
function tending to zero at infinity, hence it is identically zero.  This proves
\eqref{eq:central-identity}.

\subsection{The \texorpdfstring{$y\to1$}{y to 1} identity}

Set $y=1$ in \eqref{eq:central-identity}.  Since $U_n(1)=n+1$,
\begin{equation}
  U_{N-d-1}(1)+U_{d-1}(1)=(N-d)+d=N,
  \qquad
  U_{N-1}(1)=N.
\end{equation}
Thus
\begin{equation}
  \boxed{
  \frac1N\sum_{j=1}^{N-1}
  \frac{\bigl[1-(-1)^j\bigr]\sin(d\theta_j)\sin\theta_j}
  {1-\cos\theta_j}=1.}
  \label{eq:identity-y1}
\end{equation}

\subsection{The \texorpdfstring{$y=\cos(2\pi r/N)$}{y = cos(2pi r/N)} identity}

For $r=0,1,\ldots,N-1$, set
\begin{equation}
  y=\cos\frac{2\pi r}{N}.
\end{equation}
If this value is not a pole, the substitution is direct.  If it is a zero of $U_{N-1}$, the value
is the removable limit already determined by the residue calculation above.  Using the
trigonometric form gives
\begin{equation}
  \boxed{
  \frac1N\sum_{j=1}^{N-1}
  \frac{\bigl[1-(-1)^j\bigr]\sin(d\theta_j)\sin\theta_j}
  {\cos(2\pi r/N)-\cos\theta_j}
  =\cos\frac{2\pi r d}{N}.}
  \label{eq:identity-y-even}
\end{equation}
The case $r=0$ is the same as \eqref{eq:identity-y1} because
$\cos(0)=1$ and $\cos(0\cdot d)=1$.

\subsection{The reciprocal Fourier identity}

For $j=1,\ldots,N-1$ and $0\leq d\leq N$, the second identity needed in the absorbing-propagator
calculation is
\begin{equation}
  \boxed{
  \frac1N\sum_{r=1}^{N-1}
  \frac{\cos(2\pi r d/N)}{\cos\theta_j-\cos(2\pi r/N)}
  =\frac{1}{N(1-\cos\theta_j)}-\frac{\sin(d\theta_j)}{\sin\theta_j}.}
  \label{eq:reciprocal-identity}
\end{equation}
To see why this is the correct form, include temporarily the omitted $r=0$ term.  Since
\begin{equation}
  \frac{1}{N}\frac{1}{\cos\theta_j-1}
  =-\frac{1}{N(1-\cos\theta_j)},
\end{equation}
\eqref{eq:reciprocal-identity} is equivalent to
\begin{equation}
  \boxed{
  \frac1N\sum_{r=0}^{N-1}
  \frac{\cos(2\pi r d/N)}{\cos\theta_j-\cos(2\pi r/N)}
  =-\frac{\sin(d\theta_j)}{\sin\theta_j}.}
  \label{eq:reciprocal-identity-full}
\end{equation}
The right-hand side satisfies the homogeneous second-order recurrence
\begin{equation}
  g_{d+1}-2\cos\theta_j\,g_d+g_{d-1}=0,
  \qquad
  g_d=-\frac{\sin(d\theta_j)}{\sin\theta_j},
  \label{eq:g-recurrence}
\end{equation}
for $1\leq d\leq N-1$.  The Fourier sum on the left-hand side satisfies the same recurrence away
from $d=0$ because multiplication by $2\cos(2\pi r/N)$ shifts the Fourier modes.  The boundary
condition is also consistent because $N\theta_j=\pi j$ implies
\begin{equation}
  g_{N-d}=-\frac{\sin(N\theta_j-d\theta_j)}{\sin\theta_j}
  =-\frac{-(-1)^j\sin(d\theta_j)}{\sin\theta_j}
  =(-1)^j\frac{\sin(d\theta_j)}{\sin\theta_j}.
\end{equation}
In the only use below, this identity is multiplied by $1-(-1)^j$, so only odd $j$ contributes,
and for odd $j$ the symmetry becomes $g_{N-d}=g_d$, as required by the cosine Fourier sum.  This
establishes \eqref{eq:reciprocal-identity} in the sector used by the calculation.  Equivalently,
one may derive it from the partial-fraction expansion of $[T_N(x)-1]^{-1}$.

\section{Absorbing ring without shortcut}
\label{sec:absorbing-ring}

\subsection{Probability-conserving ring propagator}

The probability-conserving ring propagator is
\begin{equation}
  Q_{n_0}(n,t)=\frac1N\sum_{r=0}^{N-1}
  \cos\frac{2\pi r(n-n_0)}{N}\,
  \gamma_r^t,
  \label{eq:Q-full}
\end{equation}
where $\gamma_0=1$ and $\gamma_r$ for $r\geq1$ is given by
\eqref{eq:gamma-definition}.  Separating the zero mode gives
\begin{equation}
  Q_{n_0}(n,t)=\frac1N+\frac1N\sum_{r=1}^{N-1}
  \cos\frac{2\pi r(n-n_0)}{N}\,\gamma_r^t.
  \label{eq:Q-separated}
\end{equation}

\subsection{First-passage probability to a single absorbing site}

Let the absorbing site be $m$ and define
\begin{equation}
  d_0=\abs{n_0-m},\qquad 0<d_0<N.
\end{equation}
The first-passage probability to $m$ on the ring without shortcut is
\begin{equation}
  \boxed{
  \mathcal F_{n_0}(m,t)=
  \frac{q}{N}\sum_{j=1}^{N-1}
  \bigl[1-(-1)^j\bigr]\sin(d_0\theta_j)\sin\theta_j\,
  \lambda_j^{t-1},
  \qquad t\geq1,}
  \label{eq:F0-time}
\end{equation}
where
\begin{equation}
  \lambda_j=1-q+q\cos\theta_j.
\end{equation}
If $N=2L$, only the odd $j$ values contribute, and \eqref{eq:F0-time} becomes
\begin{equation}
  \mathcal F_{n_0}(m,t)=\sum_{\ell=1}^{L}f_\ell(n_0,m)\alpha_\ell^{t-1},
  \label{eq:F0-time-odd}
\end{equation}
with
\begin{equation}
  \boxed{
  f_\ell(n_0,m)=\frac{2q}{N}
  \sin\frac{(2\ell-1)\pi\abs{n_0-m}}{N}
  \sin\frac{(2\ell-1)\pi}{N}.}
  \label{eq:f-ell}
\end{equation}

The generating function follows directly from \eqref{eq:denominator-y}:
\begin{align}
  \wt{\mathcal F}_{n_0}(m,z)
  &=z\frac{q}{N}\sum_{j=1}^{N-1}
  \frac{\bigl[1-(-1)^j\bigr]\sin(d_0\theta_j)\sin\theta_j}
  {1-z\lambda_j} \notag\\
  &=\frac1N\sum_{j=1}^{N-1}
  \frac{\bigl[1-(-1)^j\bigr]\sin(d_0\theta_j)\sin\theta_j}
  {y-\cos\theta_j}\notag\\
  &=\boxed{\frac{U_{N-d_0-1}(y)+U_{d_0-1}(y)}{U_{N-1}(y)}}.
  \label{eq:F0-gen}
\end{align}
Setting $z=1$ gives $y=1$, and \eqref{eq:identity-y1} gives
\begin{equation}
  \wt{\mathcal F}_{n_0}(m,1)=1.
\end{equation}
Thus the absorbing target is eventually reached with probability one.

\subsection{Survival probability}

The survival probability up to and including time $t$ is
\begin{equation}
  R_{n_0}(t)=1-\sum_{s=0}^{t}\mathcal F_{n_0}(m,s).
\end{equation}
Using \eqref{eq:F0-time}, $\mathcal F_{n_0}(m,0)=0$, and the finite geometric sum
\begin{equation}
  \sum_{s=1}^{t}\lambda_j^{s-1}=\frac{1-\lambda_j^t}{1-\lambda_j}
  =\frac{1-\lambda_j^t}{q(1-\cos\theta_j)},
\end{equation}
we get
\begin{align}
  \sum_{s=1}^{t}\mathcal F_{n_0}(m,s)
  &=\frac1N\sum_{j=1}^{N-1}
  \frac{\bigl[1-(-1)^j\bigr]\sin(d_0\theta_j)\sin\theta_j}
  {1-\cos\theta_j}
  \bigl(1-\lambda_j^t\bigr)\notag\\
  &=1-\frac1N\sum_{j=1}^{N-1}
  \frac{\bigl[1-(-1)^j\bigr]\sin(d_0\theta_j)\sin\theta_j}
  {1-\cos\theta_j}\lambda_j^t,
\end{align}
where the last step used \eqref{eq:identity-y1}.  Hence
\begin{equation}
  \boxed{
  R_{n_0}(t)=\frac1N\sum_{j=1}^{N-1}
  \frac{\bigl[1-(-1)^j\bigr]\sin(d_0\theta_j)\sin\theta_j}
  {1-\cos\theta_j}\lambda_j^t.}
  \label{eq:survival}
\end{equation}

\subsection{Renewal equation and the killed propagator}

For $t\geq1$, the killed propagator satisfies the renewal relation
\begin{equation}
  W_{n_0}(n,t)=Q_{n_0}(n,t)
  -\sum_{s=1}^{t}\mathcal F_{n_0}(m,s)Q_m(n,t-s).
  \label{eq:renewal}
\end{equation}
This subtracts all paths that first hit $m$ at time $s$ and then, in the unrestricted ring,
continue from $m$ to $n$ in the remaining $t-s$ steps.

Insert \eqref{eq:Q-separated} and \eqref{eq:F0-time}.  The zero Fourier mode of $Q_m$ gives
\begin{equation}
  \frac1N\sum_{s=1}^{t}\mathcal F_{n_0}(m,s)
  =\frac1N\bigl[1-R_{n_0}(t)\bigr].
\end{equation}
This cancels the $1/N$ term in $Q_{n_0}(n,t)$ and leaves $R_{n_0}(t)/N$, which is precisely the
second line in \eqref{eq:W-expanded}.  For the nonzero Fourier modes, we need
\begin{align}
  \sum_{s=1}^{t}\lambda_j^{s-1}\gamma_r^{t-s}
  &=\gamma_r^{t-1}+\lambda_j\gamma_r^{t-2}+\cdots+\lambda_j^{t-1}\notag\\
  &=\frac{\gamma_r^t-\lambda_j^t}{\gamma_r-\lambda_j}\notag\\
  &=\frac{\gamma_r^t-\lambda_j^t}
  {q\bigl[\cos(2\pi r/N)-\cos\theta_j\bigr]}.
  \label{eq:time-convolution}
\end{align}
Because this convolution is subtracted in \eqref{eq:renewal}, the denominator can be rewritten as
\[
  \cos\theta_j-\cos(2\pi r/N),
\]
giving the fully expanded expression
\begin{align}
  W_{n_0}(n,t)
  &=\frac1N\sum_{r=1}^{N-1}
  \cos\frac{2\pi r(n-n_0)}{N}\,\gamma_r^t\notag\\
  &\quad+\frac{1}{N^2}\sum_{j=1}^{N-1}
  \frac{\bigl[1-(-1)^j\bigr]\sin(d_0\theta_j)\sin\theta_j}
  {1-\cos\theta_j}\lambda_j^t\notag\\
  &\quad+\frac{1}{N^2}\sum_{j=1}^{N-1}\sum_{r=1}^{N-1}
  \frac{\bigl[1-(-1)^j\bigr]\sin(d_0\theta_j)\sin\theta_j
  \cos\frac{2\pi r(n-m)}{N}}
  {\cos\theta_j-\cos(2\pi r/N)}
  \bigl(\gamma_r^t-\lambda_j^t\bigr).
  \label{eq:W-expanded}
\end{align}
The sign of the last bracket is important.  It must be $\gamma_r^t-\lambda_j^t$ when the
denominator is $\cos\theta_j-\cos(2\pi r/N)$.

Now reduce the double sum.  The part proportional to $\gamma_r^t$ uses
\eqref{eq:identity-y-even} with $d=d_0$:
\begin{align}
  &\sum_{j=1}^{N-1}
  \frac{\bigl[1-(-1)^j\bigr]\sin(d_0\theta_j)\sin\theta_j}
  {\cos\theta_j-\cos(2\pi r/N)}\notag\\
  &\qquad =-N\cos\frac{2\pi r d_0}{N}.
\end{align}
Therefore the $\gamma_r^t$ part of the double sum is
\begin{equation}
  -N\sum_{r=1}^{N-1}
  \cos\frac{2\pi r d_0}{N}
  \cos\frac{2\pi r(n-m)}{N}\,\gamma_r^t.
  \label{eq:double-gamma-part}
\end{equation}
The part proportional to $-\lambda_j^t$ uses \eqref{eq:reciprocal-identity} with
$d=\abs{n-m}$:
\begin{align}
  &-\sum_{j=1}^{N-1}\bigl[1-(-1)^j\bigr]\sin(d_0\theta_j)\sin\theta_j\,
  \lambda_j^t
  \sum_{r=1}^{N-1}\frac{\cos(2\pi r(n-m)/N)}
  {\cos\theta_j-\cos(2\pi r/N)}\notag\\
  &=-\sum_{j=1}^{N-1}\bigl[1-(-1)^j\bigr]\sin(d_0\theta_j)\sin\theta_j\,
  \lambda_j^t
  \left[\frac{1}{1-\cos\theta_j}
  -N\frac{\sin(d\theta_j)}{\sin\theta_j}\right]\notag\\
  &=-\sum_{j=1}^{N-1}\bigl[1-(-1)^j\bigr]
  \left[\frac{\sin(d_0\theta_j)\sin\theta_j}{1-\cos\theta_j}
  -N\sin(d_0\theta_j)\sin(d\theta_j)\right]\lambda_j^t.
  \label{eq:double-lambda-part}
\end{align}
After substituting \eqref{eq:double-gamma-part} and \eqref{eq:double-lambda-part} into
\eqref{eq:W-expanded}, the survival term cancels the first term inside the square brackets of
\eqref{eq:double-lambda-part}.  The final killed propagator is
\begin{equation}
  \boxed{\begin{aligned}
  W_{n_0}(n,t)
  &=\frac1N\sum_{r=1}^{N-1}
  \left[
  \cos\frac{2\pi r(n-n_0)}{N}
  -\cos\frac{2\pi r(n-m)}{N}\cos\frac{2\pi r(n_0-m)}{N}
  \right]\gamma_r^t\\
  &\quad+\frac1N\sum_{j=1}^{N-1}\bigl[1-(-1)^j\bigr]
  \sin\frac{\pi j\abs{n_0-m}}{N}
  \sin\frac{\pi j\abs{n-m}}{N}\,\lambda_j^t.
  \end{aligned}}
  \label{eq:W-final}
\end{equation}

\subsection{Boundary and initial-condition checks}

If $n=m$, then the sine factor in the second line of \eqref{eq:W-final} is zero and the first-line
bracket becomes
\begin{equation}
  \cos\frac{2\pi r(m-n_0)}{N}
  -1\cdot\cos\frac{2\pi r(n_0-m)}{N}=0.
\end{equation}
So $W_{n_0}(m,t)=0$, as required for a killed propagator.  Similarly, if $n_0=m$, both lines
vanish for transient $n$.

At $t=0$, for transient $n,n_0\neq m$, \eqref{eq:W-final} gives the identity
$W_{n_0}(n,0)=\delta_{n,n_0}$.  Here are the details.  Let
\begin{equation}
  d=\abs{n-m},\qquad d_0=\abs{n_0-m}.
\end{equation}
The basic Fourier identity is
\begin{equation}
  \frac1N\sum_{r=1}^{N-1}\cos\frac{2\pi r(n-n_0)}{N}
  =\delta_{n,n_0}-\frac1N.
\end{equation}
The product-to-sum formula gives
\begin{align}
  \sum_{r=1}^{N-1}\cos\frac{2\pi r d}{N}\cos\frac{2\pi r d_0}{N}
  &=\frac12\sum_{r=1}^{N-1}
  \left[\cos\frac{2\pi r(d-d_0)}{N}+\cos\frac{2\pi r(d+d_0)}{N}\right]\notag\\
  &=\frac{N}{2}\bigl(\delta_{d,d_0}+\delta_{d,N-d_0}\bigr)-1,
  \label{eq:cos-orthogonality}
\end{align}
where the Kronecker deltas are interpreted modulo the ring.  The odd sine sum is
\begin{align}
  \sum_{j=1}^{N-1}\bigl[1-(-1)^j\bigr]\sin(d\theta_j)\sin(d_0\theta_j)
  &=2\sum_{\substack{1\leq j\leq N-1\\j\;\mathrm{odd}}}\sin(d\theta_j)\sin(d_0\theta_j)\notag\\
  &=\frac{N}{2}\bigl(\delta_{d,d_0}+\delta_{d,N-d_0}\bigr).
  \label{eq:odd-sine-orthogonality}
\end{align}
Substituting \eqref{eq:cos-orthogonality} and \eqref{eq:odd-sine-orthogonality} into
\eqref{eq:W-final} at $t=0$ cancels the distance-only terms and leaves
$\delta_{n,n_0}$.

\subsection{Generating-function form}

Define
\begin{align}
  G^{(1)}_r(n_0,n,m)&=\frac{1}{N}
  \left[
  \cos\frac{2\pi r(n-n_0)}{N}
  -\cos\frac{2\pi r(n-m)}{N}\cos\frac{2\pi r(n_0-m)}{N}
  \right],\\
  G^{(2)}_j(n_0,n,m)&=\frac{1}{N}\bigl[1-(-1)^j\bigr]
  \sin\frac{\pi j\abs{n_0-m}}{N}
  \sin\frac{\pi j\abs{n-m}}{N}.
\end{align}
Then
\begin{equation}
  W_{n_0}(n,t)=\sum_{r=1}^{N-1}G^{(1)}_r(n_0,n,m)\gamma_r^t
  +\sum_{j=1}^{N-1}G^{(2)}_j(n_0,n,m)\lambda_j^t.
  \label{eq:W-compact-full}
\end{equation}
If $N=2L$, the second sum can be written over odd modes only:
\begin{equation}
  W_{n_0}(n,t)=\sum_{r=1}^{N-1}G^{(1)}_r(n_0,n,m)\gamma_r^t
  +\sum_{\ell=1}^{L}g^{(2)}_\ell(n_0,n,m)\alpha_\ell^t,
  \label{eq:W-compact-odd}
\end{equation}
where
\begin{equation}
  g^{(2)}_\ell(n_0,n,m)=\frac{2}{N}
  \sin\frac{(2\ell-1)\pi\abs{n_0-m}}{N}
  \sin\frac{(2\ell-1)\pi\abs{n-m}}{N}.
\end{equation}
The generating function is therefore
\begin{equation}
  \wt W_{n_0}(n,z)=\sum_{r=1}^{N-1}\frac{G^{(1)}_r(n_0,n,m)}{1-z\gamma_r}
  +\sum_{\ell=1}^{L}\frac{g^{(2)}_\ell(n_0,n,m)}{1-z\alpha_\ell}.
  \label{eq:W-gen-compact}
\end{equation}

\subsection{Special symmetric return propagator}
\label{sec:WLL}

Let $N=2L$, $m=0$, and $n=n_0=L$.  Then the first line of \eqref{eq:W-final} vanishes because
\begin{equation}
  \cos\frac{2\pi r(L-L)}{N}=1,
  \qquad
  \cos\frac{2\pi rL}{N}\cos\frac{2\pi rL}{N}=\cos^2(\pi r)=1.
\end{equation}
The second line becomes
\begin{equation}
  W_L(L,t)=\frac1N\sum_{j=1}^{N-1}\bigl[1-(-1)^j\bigr]\sin^2\frac{\pi j}{2}\,\lambda_j^t.
\end{equation}
Only odd $j=2\ell-1$ contribute.  For odd $j$, $1-(-1)^j=2$ and
$\sin^2(\pi j/2)=1$, so
\begin{equation}
  W_L(L,t)=\frac{2}{N}\sum_{\ell=1}^{L}\alpha_\ell^t.
  \label{eq:WLL-time}
\end{equation}
Hence
\begin{align}
  \wt W_L(L,z)
  &=\frac{2}{N}\sum_{\ell=1}^{L}\frac{1}{1-z\alpha_\ell}
  =\frac{2}{N}\frac{1}{zq}
  \sum_{\ell=1}^{L}\frac{1}{y-\cos\frac{(2\ell-1)\pi}{2L}}.\label{eq:WLL-gen-step1}
\end{align}
The zeros of $T_L(y)$ are exactly $\cos[(2\ell-1)\pi/(2L)]$, so
\begin{equation}
  \frac{T_L'(y)}{T_L(y)}=\sum_{\ell=1}^{L}
  \frac{1}{y-\cos\frac{(2\ell-1)\pi}{2L}}.
\end{equation}
Since $T_L'(y)=L U_{L-1}(y)$ and $N=2L$, \eqref{eq:WLL-gen-step1} gives
\begin{equation}
  \wt W_L(L,z)=\frac{2}{2L}\frac1{zq}\frac{L U_{L-1}(y)}{T_L(y)}
  =\frac{U_{L-1}(y)}{zq\,T_L(y)}.
\end{equation}
Using $U_{2L-1}(y)=2T_L(y)U_{L-1}(y)$, this is also
\begin{equation}
  \boxed{
  \wt W_L(L,z)=\frac{2U_{L-1}^2(y)}{zq\,U_{2L-1}(y)}.}
  \label{eq:WLL-gen}
\end{equation}
At $z=1$, $y=1$, so
\begin{equation}
  \wt W_L(L,1)=\frac{U_{L-1}(1)}{qT_L(1)}=\frac{L}{q}=\frac{N}{2q}.
\end{equation}

\section{Shortcut into the absorbing site}
\label{sec:shortcut}

\subsection{Transient matrix derivation of the sign}

Let $K_0$ be the transient transition matrix for the ring with absorbing target $v$, but without
the shortcut.  Its resolvent is
\begin{equation}
  R_0(z)=(I-zK_0)^{-1},
  \qquad
  [R_0(z)]_{n,n_0}=\wt W_{n_0}(n,z).
\end{equation}
A shortcut from $u$ to the absorbing target $v$ that replaces probability
$\lambda=\beta(1-q)$ of the stay-put move at $u$ removes that probability from the transient
state $u$.  Therefore, in the transient subspace,
\begin{equation}
  K=K_0-\lambda |u\rangle\langle u|.
  \label{eq:K-with-killing}
\end{equation}
The resolvent with the shortcut is
\begin{equation}
  R(z)=(I-zK)^{-1}=\bigl[I-zK_0+z\lambda |u\rangle\langle u|\bigr]^{-1}.
\end{equation}
Set
\begin{equation}
  A=I-zK_0.
\end{equation}
Then
\begin{equation}
  R(z)=\bigl[A+z\lambda |u\rangle\langle u|\bigr]^{-1}.
\end{equation}
Use the Sherman--Morrison identity
\begin{equation}
  (A+|p\rangle\langle q|)^{-1}
  =A^{-1}-\frac{A^{-1}|p\rangle\langle q|A^{-1}}
  {1+\langle q|A^{-1}|p\rangle}.
\end{equation}
Here $|p\rangle=z\lambda |u\rangle$ and $\langle q|=\langle u|$.  Thus
\begin{align}
  R(z)&=R_0(z)-\frac{z\lambda R_0(z)|u\rangle\langle u|R_0(z)}
  {1+z\lambda\langle u|R_0(z)|u\rangle}.
\end{align}
Taking the $(n,n_0)$ matrix element gives the corrected propagator
\begin{equation}
  \boxed{
  \wt S_{n_0}(n,z)=\wt W_{n_0}(n,z)-
  \frac{z\lambda\,\wt W_{n_0}(u,z)\wt W_u(n,z)}
  {1+z\lambda\,\wt W_u(u,z)}.}
  \label{eq:S-corrected}
\end{equation}
This is the same as \eqref{eq:main-corrected-resolvent}.  The sign is not a convention: it is fixed
by the fact that the target state $v$ is absorbing and therefore absent from the transient
subspace.

\subsection{First-passage generating function}

Sum \eqref{eq:S-corrected} over all transient sites $n\neq v$.  For the baseline killed walk,
\begin{equation}
  \sum_{n\neq v}\wt W_r(n,z)=\frac{1-\wt F^{(0)}_r(v,z)}{1-z}.
  \label{eq:survival-gen-relation}
\end{equation}
Therefore
\begin{align}
  \sum_{n\neq v}\wt S_{n_0}(n,z)
  &=\frac{1-\wt F^{(0)}_{n_0}(v,z)}{1-z}
  -\frac{z\lambda\,\wt W_{n_0}(u,z)}{1+z\lambda\wt W_u(u,z)}
  \frac{1-\wt F^{(0)}_u(v,z)}{1-z}.
\end{align}
But the left-hand side is also $(1-\wt F_{n_0}(v,z))/(1-z)$.  Multiplying by $1-z$ and rearranging
gives
\begin{equation}
  \boxed{
  \wt F_{n_0}(v,z)=\wt F^{(0)}_{n_0}(v,z)+
  \frac{z\lambda\,\wt W_{n_0}(u,z)\bigl[1-\wt F^{(0)}_u(v,z)\bigr]}
  {1+z\lambda\,\wt W_u(u,z)}.}
  \label{eq:F-corrected-shortcut}
\end{equation}

\subsection{Why the sign is also forced physically}

For $0\leq z<1$, all coefficients of $\wt W_{n_0}(u,z)$ and
$1-\wt F^{(0)}_u(v,z)$ are non-negative, and $1+z\lambda\wt W_u(u,z)>0$.  Hence the shortcut term
in \eqref{eq:F-corrected-shortcut} is non-negative.  Thus adding a shortcut into the absorbing
site increases the first-passage generating function, as it should.

By contrast, a formula with
\begin{equation}
  +\frac{z\lambda\wt W_{n_0}(u,z)\wt W_u(n,z)}{1-z\lambda\wt W_u(u,z)}
\end{equation}
in the propagator would increase transient survival and decrease the hitting generating function.
It can also produce a positive real singularity when
\begin{equation}
  z\lambda\wt W_u(u,z)=1.
\end{equation}
In the symmetric case, $\wt W_L(L,1)=N/(2q)$, so the denominator at $z=1$ would be
\begin{equation}
  1-\frac{\beta(1-q)N}{2q},
\end{equation}
which is exactly the source of the original threshold.  Such a singularity is an artefact of the
wrong sign; it is not present in the corrected stochastic shortcut-to-target model.

\section{Symmetric shortcut: exact corrected first-passage generating function}
\label{sec:symmetric-shortcut}

We now set
\begin{equation}
  N=2L,
  \qquad v=0,
  \qquad u=L.
\end{equation}
Let $F_\rho(z)$ be the first-passage generating function to $0$ from a site at distance
$\rho$ from $0$ along either half-ring, where $0\leq\rho\leq L$.  By symmetry,
$F_{L+k}=F_{L-k}$.

\subsection{Difference equation}

For $1\leq\rho\leq L-1$, the first step gives
\begin{equation}
  F_\rho(z)=z\left[(1-q)F_\rho(z)+\frac q2F_{\rho-1}(z)+\frac q2F_{\rho+1}(z)\right].
\end{equation}
Move all terms to one side and divide by $zq/2$:
\begin{equation}
  F_{\rho+1}(z)-2yF_\rho(z)+F_{\rho-1}(z)=0,
  \qquad
  y=1+\frac1q\left(\frac1z-1\right).
  \label{eq:interior-recurrence}
\end{equation}
The absorbing boundary is
\begin{equation}
  F_0(z)=1.
\end{equation}
At $\rho=L$, the shortcut replaces probability $\lambda=\beta(1-q)$ of the stay-put probability by
an immediate jump to the target.  Therefore
\begin{equation}
  F_L(z)=z\left[\lambda+(1-q-\lambda)F_L(z)+qF_{L-1}(z)\right].
\end{equation}
Using $1-q-\lambda=(1-\beta)(1-q)$ and writing
\begin{equation}
  h=\frac{\lambda}{q}=\frac{\beta(1-q)}{q},
  \qquad a=\frac1h=\frac{q}{\beta(1-q)},
\end{equation}
the boundary equation becomes
\begin{equation}
  \boxed{(y+h)F_L(z)-F_{L-1}(z)=h.}
  \label{eq:L-boundary}
\end{equation}

\subsection{Solving the recurrence}

The recurrence \eqref{eq:interior-recurrence} is the Chebyshev recurrence.  The solution satisfying
$F_0=1$ can be written as
\begin{equation}
  F_\rho=T_\rho(y)+C\,U_{\rho-1}(y),
  \label{eq:F-rho-general}
\end{equation}
because $T_0=1$ and $U_{-1}=0$.  Hence
\begin{equation}
  F_L=T_L+C U_{L-1},
  \qquad
  F_{L-1}=T_{L-1}+C U_{L-2}.
\end{equation}
Substitute these into \eqref{eq:L-boundary}:
\begin{align}
  h&=(y+h)(T_L+C U_{L-1})-(T_{L-1}+C U_{L-2})\notag\\
  &=\bigl(yT_L-T_{L-1}\bigr)
  +C\bigl(yU_{L-1}-U_{L-2}\bigr)
  +hT_L+hC U_{L-1}.
  \label{eq:C-equation-step1}
\end{align}
The needed Chebyshev identities are
\begin{equation}
  yU_{L-1}-U_{L-2}=T_L,
  \qquad
  yT_L-T_{L-1}=(y^2-1)U_{L-1}.
  \label{eq:cheb-identities-boundary}
\end{equation}
Using them in \eqref{eq:C-equation-step1} gives
\begin{equation}
  h=(y^2-1)U_{L-1}+C T_L+hT_L+hC U_{L-1}.
\end{equation}
Therefore
\begin{equation}
  C\bigl(T_L+hU_{L-1}\bigr)=h(1-T_L)-(y^2-1)U_{L-1},
\end{equation}
and
\begin{equation}
  C=\frac{h(1-T_L)-(y^2-1)U_{L-1}}{T_L+hU_{L-1}}.
  \label{eq:C-value}
\end{equation}
Substitution into \eqref{eq:F-rho-general} gives one correct form of the answer:
\begin{equation}
  F_\rho=T_\rho+U_{\rho-1}
  \frac{h(1-T_L)-(y^2-1)U_{L-1}}{T_L+hU_{L-1}}.
  \label{eq:F-rho-form1}
\end{equation}
Multiplying numerator and denominator by $a=1/h$ gives
\begin{equation}
  F_\rho=\frac{aT_\rho T_L+T_\rho U_{L-1}+U_{\rho-1}(1-T_L)
  -a(y^2-1)U_{\rho-1}U_{L-1}}{aT_L+U_{L-1}}.
  \label{eq:F-rho-form2}
\end{equation}
Now use the two product identities
\begin{align}
  T_\rho T_L-(y^2-1)U_{\rho-1}U_{L-1}&=T_{L-\rho},
  \label{eq:product-id1}\\
  T_\rho U_{L-1}+U_{\rho-1}(1-T_L)&=U_{L-\rho-1}+U_{\rho-1}.
  \label{eq:product-id2}
\end{align}
For example, if $y=\cos\theta$, then \eqref{eq:product-id1} reads
\begin{equation}
  \cos(\rho\theta)\cos(L\theta)+\sin(\rho\theta)\sin(L\theta)
  =\cos((L-\rho)\theta),
\end{equation}
and \eqref{eq:product-id2} follows from
\begin{equation}
  \cos(\rho\theta)\frac{\sin(L\theta)}{\sin\theta}
  +\frac{\sin(\rho\theta)}{\sin\theta}\bigl[1-\cos(L\theta)\bigr]
  =\frac{\sin((L-\rho)\theta)+\sin(\rho\theta)}{\sin\theta}.
\end{equation}
Thus \eqref{eq:F-rho-form2} becomes
\begin{equation}
  \boxed{
  \wt F_\rho(z)=
  \frac{aT_{L-\rho}(y)+U_{\rho-1}(y)+U_{L-\rho-1}(y)}
  {aT_L(y)+U_{L-1}(y)}.}
  \label{eq:F-rho-final}
\end{equation}
This is the same as \eqref{eq:main-symmetric-F}.

\subsection{Checks on the symmetric formula}

First, $\rho=0$ gives
\begin{equation}
  \wt F_0(z)=\frac{aT_L+U_{-1}+U_{L-1}}{aT_L+U_{L-1}}=1,
\end{equation}
as required for a walker already at the absorbing site.  Second, in the no-shortcut limit
$\beta\to0$, we have $a\to\infty$, so
\begin{equation}
  \wt F_\rho(z)\to\frac{T_{L-\rho}(y)}{T_L(y)}.
\end{equation}
This agrees with the absorbing-ring formula \eqref{eq:F0-gen}.  Indeed, for $N=2L$,
\begin{equation}
  \frac{U_{2L-\rho-1}+U_{\rho-1}}{U_{2L-1} }
  =\frac{T_{L-\rho}}{T_L},
\end{equation}
which follows from the same trigonometric product identities.

Third, when $\rho=L$,
\begin{equation}
  \wt F_L(z)=\frac{a+U_{L-1}(y)}{aT_L(y)+U_{L-1}(y)}.
\end{equation}
At $z=1$, $y=1$, so $T_L(1)=1$ and $U_{L-1}(1)=L$, giving
\begin{equation}
  \wt F_L(1)=\frac{a+L}{a+L}=1.
\end{equation}
Thus absorption remains certain.

\section{Poles, tail decay, and the corrected dependence on \texorpdfstring{$\beta$}{beta}}
\label{sec:poles-tail}

\subsection{Pole equation}

Write $s=1/z$, so
\begin{equation}
  y_s=1+\frac{s-1}{q}.
\end{equation}
The poles of \eqref{eq:F-rho-final} are roots of
\begin{equation}
  D(y_s)=aT_L(y_s)+U_{L-1}(y_s)=0.
  \label{eq:D-y}
\end{equation}
If $y_s=\cos\theta$, then
\begin{equation}
  T_L(y_s)=\cos(L\theta),
  \qquad
  U_{L-1}(y_s)=\frac{\sin(L\theta)}{\sin\theta}.
\end{equation}
Therefore \eqref{eq:D-y} becomes
\begin{equation}
  a\cos(L\theta)+\frac{\sin(L\theta)}{\sin\theta}=0.
\end{equation}
Multiplying by $\sin\theta$ and dividing by $\cos(L\theta)$ away from the asymptotes gives
\begin{equation}
  \boxed{\tan(L\theta)=-a\sin\theta.}
  \label{eq:theta-root}
\end{equation}
The corresponding pole value is
\begin{equation}
  \boxed{s=1-q+q\cos\theta.}
  \label{eq:s-theta}
\end{equation}

\subsection{Location of the largest pole}

The no-shortcut largest odd-sector eigenvalue is
\begin{equation}
  \alpha_1=1-q+q\cos\frac{\pi}{2L}=1-q+q\cos\frac{\pi}{N}.
\end{equation}
The largest even-sector nontrivial eigenvalue is
\begin{equation}
  \gamma_1=1-q+q\cos\frac{\pi}{L}=1-q+q\cos\frac{2\pi}{N}.
\end{equation}
Consider the interval
\begin{equation}
  \frac{\pi}{2L}<\theta<\frac{\pi}{L}.
\end{equation}
In this interval,
\begin{equation}
  \cos(L\theta)<0,
  \qquad
  \sin(L\theta)>0,
  \qquad
  \sin\theta>0.
\end{equation}
Hence $D(\cos\theta)=a\cos(L\theta)+\sin(L\theta)/\sin\theta$ tends to $U_{L-1}>0$ at the left
endpoint $\theta=\pi/(2L)$, while it equals $a\cos\pi=-a<0$ at the right endpoint
$\theta=\pi/L$.  By continuity there is a root $\theta_1$ in this interval.  This root gives
\begin{equation}
  \cos\frac{\pi}{L}<\cos\theta_1<\cos\frac{\pi}{2L},
\end{equation}
and therefore
\begin{equation}
  \boxed{\gamma_1<s_1<\alpha_1.}
\end{equation}
Since the transient transition matrix is non-negative and primitive for $0<q<1$, its spectral
radius is the unique positive dominant eigenvalue.  Thus $s_1$ controls the long-time tail.

\subsection{Partial fractions and the exact tail coefficient}

Let
\begin{equation}
  N_\rho(y)=aT_{L-\rho}(y)+U_{\rho-1}(y)+U_{L-\rho-1}(y),
  \qquad
  D(y)=aT_L(y)+U_{L-1}(y).
\end{equation}
Let $y_k$ be a simple root of $D(y)$ and define
\begin{equation}
  s_k=1-q+qy_k.
\end{equation}
Near $s=s_k$,
\begin{equation}
  \wt F_\rho(1/s)=\frac{N_\rho(y_s)}{D(y_s)}
  =\frac{N_\rho(y_k)}{D'(y_k)(y_s-y_k)}+\text{regular terms}.
\end{equation}
Since $y_s-y_k=(s-s_k)/q$, the residue in the $s$-plane is
\begin{equation}
  \Res_{s=s_k}\wt F_\rho(1/s)=\frac{qN_\rho(y_k)}{D'(y_k)}.
\end{equation}
The partial-fraction term corresponding to a pole $s_k$ is
\begin{equation}
  \frac{B_{\rho k}}{1-zs_k},
\end{equation}
whose $s=1/z$ form is
\begin{equation}
  \frac{B_{\rho k}}{1-s_k/s}=\frac{B_{\rho k}s}{s-s_k}.
\end{equation}
Its residue at $s=s_k$ is $B_{\rho k}s_k$.  Hence the coefficient multiplying
$s_k^t$ is
\begin{equation}
  B_{\rho k}^{(z^t)}=\frac{qN_\rho(y_k)}{s_kD'(y_k)}.
\end{equation}
Equivalently, if we write the time dependence as $s_k^{t-1}$ for $t\geq1$, then
\begin{equation}
  \boxed{
  F_\rho(t)=\sum_k B_{\rho k}s_k^{t-1},
  \qquad
  B_{\rho k}=\frac{qN_\rho(y_k)}{D'(y_k)}.}
  \label{eq:partial-fraction-B}
\end{equation}
In particular,
\begin{equation}
  \boxed{
  F_\rho(t)\sim B_{\rho1}s_1^{t-1},
  \qquad
  B_{\rho1}=\frac{qN_\rho(y_1)}{D'(y_1)}.}
  \label{eq:tail-coefficient}
\end{equation}
The decay time can be reported as
\begin{equation}
  \boxed{\tau_\beta=-\frac{1}{\log s_1}.}
  \label{eq:tau-beta}
\end{equation}

\subsection{Weak-shortcut expansion}

For small $\beta$, equivalently large
\begin{equation}
  a=\frac{q}{\beta(1-q)},
\end{equation}
the largest root lies close to
\begin{equation}
  \theta_0=\frac{\pi}{2L}.
\end{equation}
Write
\begin{equation}
  \theta_1=\theta_0+\delta,
  \qquad 0<\delta\ll1.
\end{equation}
Then
\begin{equation}
  L\theta_1=\frac{\pi}{2}+L\delta,
  \qquad
  \tan(L\theta_1)=\tan\left(\frac\pi2+L\delta\right)
  =-\cot(L\delta).
\end{equation}
For small $\delta$,
\begin{equation}
  -\cot(L\delta)=-\frac{1}{L\delta}+O(L\delta).
\end{equation}
The root equation \eqref{eq:theta-root} gives
\begin{equation}
  -\frac{1}{L\delta}+O(\delta)=-a\sin(\theta_0+\delta).
\end{equation}
To leading order,
\begin{equation}
  \delta=\frac{1}{aL\sin\theta_0}+O(a^{-2}).
\end{equation}
Now expand $s_1=1-q+q\cos(\theta_0+\delta)$:
\begin{align}
  s_1
  &=1-q+q\left[\cos\theta_0-\delta\sin\theta_0+O(\delta^2)\right]\notag\\
  &=\alpha_1-q\sin\theta_0\frac{1}{aL\sin\theta_0}+O(a^{-2})\notag\\
  &=\alpha_1-\frac{q}{aL}+O(a^{-2}).
\end{align}
Since $q/a=\beta(1-q)$ and $L=N/2$,
\begin{equation}
  \boxed{s_1=\alpha_1-\frac{2\beta(1-q)}{N}
  +O\bigl([\beta(1-q)]^2\bigr).}
\end{equation}
Thus
\begin{equation}
  \boxed{1-s_1=1-\alpha_1+\frac{2\beta(1-q)}{N}
  +O\bigl([\beta(1-q)]^2\bigr).}
\end{equation}
For large $N$,
\begin{equation}
  1-\alpha_1=q\left[1-\cos\frac{\pi}{N}\right]
  =\frac{q\pi^2}{2N^2}+O(N^{-4}),
\end{equation}
so
\begin{equation}
  \boxed{1-s_1\sim \frac{q\pi^2}{2N^2}+\frac{2\beta(1-q)}{N}.}
\end{equation}

\subsection{Strong-shortcut endpoint expansion}

When $a$ is small, the largest root approaches the zero of $U_{L-1}$ at
$\theta=\pi/L$.  Write
\begin{equation}
  \theta_1=\frac{\pi}{L}-\varepsilon,
  \qquad 0<\varepsilon\ll1.
\end{equation}
Then
\begin{equation}
  \tan(L\theta_1)=\tan(\pi-L\varepsilon)=-L\varepsilon+O(\varepsilon^3),
\end{equation}
and
\begin{equation}
  -a\sin\theta_1=-a\sin\frac{\pi}{L}+O(a\varepsilon).
\end{equation}
Thus
\begin{equation}
  \varepsilon=\frac{a}{L}\sin\frac{\pi}{L}+O(a^2).
\end{equation}
Consequently
\begin{align}
  s_1&=1-q+q\cos\left(\frac{\pi}{L}-\varepsilon\right)\notag\\
  &=\gamma_1+q\varepsilon\sin\frac{\pi}{L}+O(\varepsilon^2)\notag\\
  &=\gamma_1+\frac{qa}{L}\sin^2\frac{\pi}{L}+O(a^2).
\end{align}
This shows explicitly how the dominant pole moves from $\alpha_1$ toward $\gamma_1$ as the
shortcut strength increases.

\section{Consequences for the bimodality discussion}
\label{sec:bimodality}

The existence of a second peak is a finite-time shape question.  It is not determined solely by
which pole is asymptotically largest.  In the corrected formula, if
\begin{equation}
  F_\rho(t)=\sum_k B_{\rho k}s_k^{t-1},
\end{equation}
then local extrema are controlled exactly by the finite difference
\begin{equation}
  F_\rho(t+1)-F_\rho(t)=\sum_k B_{\rho k}(s_k-1)s_k^{t-1}.
  \label{eq:finite-difference-bimodality}
\end{equation}
A second peak at time $t_*$ requires
\begin{equation}
  F_\rho(t_*)-F_\rho(t_*-1)>0,
  \qquad
  F_\rho(t_*+1)-F_\rho(t_*)<0.
\end{equation}
These inequalities can be checked exactly from \eqref{eq:finite-difference-bimodality}.  They do
not reduce to the original threshold $\beta\leq2q/[(1-q)N]$.

What can be stated analytically and cleanly is the decay of the right tail of a second peak if it
is present:
\begin{equation}
  \boxed{F_\rho(t)\sim B_{\rho1}s_1^{t-1},
  \qquad \frac{F_\rho(t+1)}{F_\rho(t)}\to s_1.}
\end{equation}
Thus the referee-facing corrected decay dependence is
\begin{equation}
  \boxed{\tau_\beta^{-1}=-\log s_1,\qquad
  aT_L\left(1+\frac{s_1-1}{q}\right)+U_{L-1}\left(1+\frac{s_1-1}{q}\right)=0.}
\end{equation}
For weak shortcuts this gives
\begin{equation}
  \boxed{\tau_\beta^{-1}
  =-\log\left[\alpha_1-\frac{2\beta(1-q)}{N}+O\bigl([\beta(1-q)]^2\bigr)\right].}
\end{equation}
Equivalently, in gap form,
\begin{equation}
  \boxed{1-s_1\sim \frac{q\pi^2}{2N^2}+\frac{2\beta(1-q)}{N}.}
\end{equation}
The shortcut contribution to the decay gap is the term
\begin{equation}
  \boxed{\frac{2\beta(1-q)}{N}.}
\end{equation}

\section{Audit notes: issues found in the original calculation}
\label{sec:audit-notes}

This section records the mathematical and typographical issues identified during the audit.

\subsection{Substantive mathematical issues}

\begin{enumerate}
\item The shortcut-to-target resolvent sign is wrong for the stochastic interpretation in which a
fraction $\beta(1-q)$ of the stay-put probability at $u$ is replaced by a jump to absorbing $v$.
The corrected transient perturbation is negative killing at $u$, giving
\eqref{eq:S-corrected}.

\item The original denominator $1-z\beta(1-q)\wt W_u(u,z)$ must be replaced by
$1+z\beta(1-q)\wt W_u(u,z)$.  This changes the pole equation from a positive horizontal-line
intersection to a negative one, or equivalently from $a-\phi=0$ to $a+\phi=0$.

\item The original first-passage correction has the wrong monotonicity: it decreases the hitting
generating function after adding a shortcut to the target.  The corrected formula
\eqref{eq:F-corrected-shortcut} increases it.

\item The original beta condition
\begin{equation}
  \beta\leq\frac{2q}{(1-q)N}
\end{equation}
comes from the wrong-sign denominator and is not a valid corrected condition for the stochastic
shortcut-to-target model.

\item The original $t\alpha_1^t$ conclusion comes from an artificial collision between an
uncorrected shortcut pole and an absorbing-ring pole.  In the corrected model the dominant pole is
simple and equals $s_1<\alpha_1$, so the tail is $B_{\rho1}s_1^{t-1}$.

\item The reduction of the absorbing-ring double sum contains a sign inconsistency in the
intermediate line: with denominator $\cos(\pi j/N)-\cos(2\pi r/N)$, the bracket must be
$\gamma_r^t-\lambda_j^t$.  The final killed propagator is nevertheless correct after this sign is
fixed.

\item The initial-condition check following the killed propagator uses incorrect individual
orthogonality claims.  The correct cosine and odd-sine sums are \eqref{eq:cos-orthogonality} and
\eqref{eq:odd-sine-orthogonality}; only their combination gives $\delta_{n,n_0}$.

\item In the long convolution formula corresponding to the original Eq. (41), the first
no-shortcut term should involve $f_\ell(n_0,v)$, not $f_\ell(n_0,u)$.  The baseline term in
$F_{n_0}(v,t)$ is the first-passage probability to $v$.

\item In the same long convolution formula, terms involving the odd-sector propagator
$g^{(2)}_r$ should carry $\alpha_r$, not $\gamma_r$.  The displayed use of $\gamma_r$ in the last
line is an index/eigenvalue mismatch.

\item The residue coefficient formula for the uncorrected $c_k$ contains bracket mismatches and an
apparently incorrect derivative factor.  Since that coefficient belongs to the wrong-sign pole
problem, it should be replaced rather than repaired.
\end{enumerate}

\subsection{Typographical and presentation issues}

\begin{enumerate}
\item The document class option \texttt{idelinks} should be \texttt{hidelinks}.
\item The source depends on local style files such as \texttt{customcolours.sty},
\texttt{customcommands.sty}, and \texttt{globalplotconfig.sty}; a portable derivation should avoid
these dependencies or include them.
\item A label of the form \verb|\label{eq:}| is not useful and can create reference problems.
\item Several expressions write $G^{(1)}_k(n_0,n,,m)$ and $G^{(2)}_k(n_0,n,,m)$ with a double
comma.
\item The text says $Q_{n_0}(n,0)=\delta_{n,n_0}$ and $Q_{n_0}(n,t)=0$ in places where it means
$W_{n_0}(n,0)$ and $W_{n_0}(n,t)$.
\item The compact odd-sector sums require $N$ even.  This assumption should be stated before using
$\sum_{\ell=1}^{N/2}$.
\item The figure caption says ``given the requirement $q>1$''; the stochastic model requires
$0<q<1$, and a nonzero shortcut probability requires $q<1$ if the shortcut probability is
proportional to $1-q$.
\item The phrase ``poles in the dominator'' should be ``poles in the denominator.''
\item The contour statement after $z=1/s$ should specify orientation and which singularities are
enclosed.  The corrected derivation avoids this ambiguity by using partial fractions directly in
$s$.
\end{enumerate}

\section{Optional numerical check for the figure parameters}

For the original figure parameters
\begin{equation}
  N=10,
  \qquad L=5,
  \qquad q=\frac23,
  \qquad \beta=\frac27,
\end{equation}
we have
\begin{equation}
  a=\frac{q}{\beta(1-q)}=7.
\end{equation}
The corrected largest root solves
\begin{equation}
  \tan(5\theta)=-7\sin\theta,
  \qquad \frac{\pi}{10}<\theta<\frac{\pi}{5}.
\end{equation}
Numerically,
\begin{equation}
  \theta_1\approx0.3865953134,
  \qquad
  s_1=1-\frac23+\frac23\cos\theta_1\approx0.9507987460.
\end{equation}
The no-shortcut largest odd eigenvalue is
\begin{equation}
  \alpha_1=1-\frac23+\frac23\cos\frac{\pi}{10}\approx0.9673710109,
\end{equation}
so indeed
\begin{equation}
  s_1<\alpha_1.
\end{equation}
The corrected tail gap is
\begin{equation}
  1-s_1\approx0.0492012540.
\end{equation}

\end{document}
